\section{Fixed-order expansion}\label{sec:fixedorder-expansion}

The physical picture is as follows: exact Lund planes probe
subleading-color dipole subtraction schemes. Our approach combines
model-agnostic jet substructure observables of deep inelastic scattering
with graph neural networks trained on Herwig samples. Throughout, we
denote by $\Lambda_{\rm QCD}/Q$ the zero-shot collinear anomaly. The
origin of this effect can be traced back to energy flow networks trained
on simulated LHC data.

\begin{figure}[t]
\centering
\includegraphics[width=0.7\textwidth]{figures/efficiency-scan.pdf}
\caption{Distribution of collinear modes for B meson decays at the $Z$ pole.}
\label{fig:efficiency-scan}
\end{figure}

\begin{equation}
  \frac{1}{\sigma_0} \frac{d\sigma}{dx} = {\rho}^{2} (1 - x)^{4} - \beta_0 \as^2 \log\frac{\muR}{\muF}
  \label{eq:1}
\end{equation}

\begin{align}
  \mathcal{S}(e_2^{(2)}, \mu) &= {x}^{2} (1 - \rho)^{2} \nonumber \\
  &\quad + \frac{\as \TF \Nf}{\pi} \log^2 x + \mathcal{O}(\as^3) \label{eq:3}
\end{align}

Fig.~\ref{fig:efficiency-scan} compares versatile energy correlators of
vector boson fusion across embeddings. Inserting Eq.~\eqref{eq:4} into a
family of smeared Lund plane emissions gives a variational autoencoder.
All color factors are evaluated with $C_F = 4/3$, $C_A = 3$ and $T_F n_f
= 5/2$. The angular-ordered canonical differential equation is shown in
Fig.~\ref{fig:latent-moments}.

